% ============================================
% Johns Hopkins Letter Template
% Uses the built-in 'letter' class
% ============================================

\documentclass[11pt,letterpaper]{letter}

\usepackage{graphicx}
\usepackage{eso-pic}
\usepackage{geometry}
\usepackage{microtype}
\usepackage[utf8]{inputenc}
\usepackage[hidelinks]{hyperref}

% ----- Page Setup -----
\geometry{left=25mm,right=25mm,top=25mm,bottom=25mm}

% ----- Logo Placement (foreground, first page only) -----
\newcommand{\PutJHULogoFG}[1]{%
  \AddToShipoutPictureFG*{%
    \AtPageUpperLeft{%
      \hspace{10mm}%
      \raisebox{-38mm}{%
        \includegraphics[width=6cm]{#1}%
      }%
    }%
  }%
}

% ----- Sender Information -----
\signature{Decory Edwards}
\address{Department of Economics \\ Johns Hopkins University \\ Baltimore, MD 21218 \\ \texttt{dedwar65@jh.edu}}

% ----- Recipient Info -----
\begin{document}
\PutJHULogoFG{Images/jhulogo.png}

\begin{letter}{Hiring Committee \\
American Academy of Arts and Sciences}

\opening{Dear Hiring Committee,}

I am writing to express my interest in the Program Officer position for American Institutions. I am a Ph.D.\ candidate in the Department of Economics at Johns Hopkins University (advisors: Christopher D.\ Carroll, Nicholas W.\ Papageorge, and Stelios Fourakis), and I expect to receive my degree in May 2026. My primary specializations are macroeconomic modeling, wealth and income dynamics, and computational economics.

My research agenda focuses on the ability of macroeconomic models to generate empirically realistic distributions of wealth and income over the life cycle. A key driver of that work is modeling heterogeneity in the rate of return on assets, and understanding how return dispersion contributes to portfolio accumulation across households.

In my job market paper, I calibrate a life cycle heterogeneous agent model to match wealth moments from the Survey of Consumer Finances by allowing households to differ in their portfolio returns. The model helps explain observed wealth concentration and return dispersion. In a related research project using the Health and Retirement Study, I examine how trust influences both income growth and asset returns. I find a hump shaped relationship for trust and returns (consistent with the literature) and characterize the distribution of the persistent component of returns to net worth. These experiences have equipped me to frame research strategies and design modeling approaches, execute econometric and simulation analysis with large micro data sets, and translate results into clear and rigorous written deliverables for both academic and practitioner audiences.

I am particularly drawn to the Academy’s work on American Institutions, especially its efforts to build non partisan and cross ideological projects that address economic and civic challenges. Initiatives such as the Commission on Reimagining the Economy reflect the type of work that connects empirical research with institutional design and democratic legitimacy. My own research includes analysis of capital income taxation and wealth taxation under heterogeneous returns, which speaks directly to questions about economic opportunity, fiscal design, and the structure of incentives within democratic systems. I am accustomed to synthesizing technical findings into accessible arguments that can inform broader policy conversations. The Program Officer role’s emphasis on convening scholars and civic leaders, drafting reports, and developing new initiatives aligns closely with both my substantive expertise in political economy and my experience producing research that bridges theory, data, and policy relevance.

I am enthusiastic about the opportunity to live and work in either Cambridge or Washington, DC and to participate fully in the hybrid schedule expected for this role, including regular in person collaboration. I value structured dialogue and interdisciplinary engagement, which are central to the Academy’s mission. I am also drawn to the prospect of a sustained multi year appointment, which would allow me to contribute meaningfully to the Academy’s projects while continuing to develop my research agenda in areas related to public finance, institutional design, and the political economy of inequality.

I believe I would be a strong fit for this role because:
\begin{enumerate}
    \item My expertise in political economy and public finance aligns with the Academy’s focus on the American economy and democratic institutions.
    \item My experience writing policy relevant research prepares me to draft reports, briefings, and materials for diverse audiences.
    \item I bring a collaborative approach and commitment to rigorous, non partisan analysis that supports cross disciplinary engagement.
\end{enumerate}

I would welcome the opportunity to contribute my research background, writing experience, and commitment to the public good to the Academy’s American Institutions program. Thank you for considering my application.

\closing{Sincerely,}

\end{letter}
\end{document}
